%% kapa: exam=1 round=37 year=2026 period=2 form=A subject=accounting subject_ko=회계학 seq=21 number=61
%% source: Q-Net 감정평가사 기출문제 (cst00309) artlSeq=5258924 file=raw/1008/5258924/2026년 제37회 감정평가사 1차 2교시.pdf
%% conversion: kordoc 4.24.51 -> tools/gen1.py (auto); review=proofread-sonnet; flags=none
\begin{question}{61}
(주)감평은 확정급여제도를 채택하고 있으며, 20×1년 초 순확정급여부채는 ￦20,000이다. (주)감평의 20×1년도 확정급여제도와 관련된 자료는 다음과 같다.
\begin{framedtext}
○ 순확정급여부채(자산) 계산시 적용하는 할인율은 연 10 \%이다. \par
○ 20×1년도 당기근무원가는 ￦90,000이고, 20×1년 말 퇴직종업원에게 ￦40,000의 현금이 사외적립자산에서 지급되었다. \par
○ 20×1년 말 사외적립자산에 ￦70,000을 현금으로 출연하였다. \par
○ 20×1년에 발생한 확정급여채무의 재측정요소(손실)는 ￦7,000이고, 사외적립자산의 재측정요소(이익)는 ￦3,800이다. \par
\end{framedtext}
20×1년 말 재무상태표에 보고될 순확정급여부채는? \par
\begin{choices}
  \item ￦40,800
  \item ￦41,200
  \item ￦42,000
  \item ￦44,800
  \item ￦45,200
\end{choices}
\end{question}
\answer{5}
